\begin{wraptable}{r}{0.48\linewidth}
  \vspace{-10pt}
  \centering
  \caption{Model uncertainties captured.}
  \resizebox{\linewidth}{!}{% 
  \label{tab:modeluncertainty}
  \begin{tabular}{l c c}
  \multirow{2}{*}{\textbf{Model}} & \textbf{Aleatoric} & \textbf{Epistemic}\\
  & \textbf{uncertainty}& \textbf{uncertainty}\\
  \hline \textit{Baseline Models} & & \\
  Deterministic NN (D)           & No & No\\
  Probabilistic NN (P)           & Yes & No \\
  Deterministic ensemble NN (DE) & No & Yes \\
  Gaussian process baseline (GP) & Homoscedastic & Yes \\
  \hline \textit{Our Model} & & \\
  Probabilistic ensemble NN (PE) & \textbf{Yes} & \textbf{Yes}
  \end{tabular}%
  }
%   \vspace{-13pt}
  \label{table:modelclasses}
\end{wraptable}
%
This section describes several ways to model the task's true (but unknown) dynamic function, including our method: an ensemble of bootstrapped probabilistic neural networks.
Whilst uncertainty-aware dynamics models have been explored in a number of prior works \citep{Gal2016,Depeweg2016}, the particular details of the implementation and design decisions in regard incorporation of uncertainty have not been rigorously analyzed empirically.
As a result, prior work has generally found that expressive parametric models, such as deep neural networks, generally do not produce model-based RL algorithms that are competitive with their model-free counterparts in terms of asymptotic performance \citep{Nagabandi2017}, and often even found that simpler time-varying linear models can outperform expressive neural network models \citep{Levine2016,gu2016continuous}. 

Any MBRL algorithm must select a class of model to predict the dynamics.
This choice is often crucial for an MBRL algorithm, as even small bias can significantly influence the quality of the corresponding controller~\citep{Atkeson1997,Abbeel2006}.
A major challenge is building a model that performs well in low and high data regimes: in the early stages of training, data is scarce, and highly expressive function approximators are liable to overfit;
In the later stages of training, data is plentiful, but for systems with complex dynamics, simple function approximators might underfit. 
While Bayesian models such as GPs perform well in low-data regimes, they do not scale favorably with dimensionality and often use kernels ill-suited for discontinuous dynamics~\citep{Calandra2016a}, which is typical of robots interacting through contacts.

In this paper, we study how expressive NNs can be incorporated into MBRL. 
To account for uncertainty, we study NNs that model two types of uncertainty. 
The first type, aleatoric uncertainty, arises from \textit{inherent stochasticities} of a system, \eg{} observation noise and process noise.
Aleatoric uncertainty can be captured by outputting the parameters of a parameterized distribution, while still training the network discriminatively. 
The second type -- epistemic uncertainty -- corresponds to \textit{subjective uncertainty} about the dynamics function, due to a lack of sufficient data to uniquely determine the underlying system exactly.
In the limit of infinite data, epistemic uncertainty should vanish, but for datasets of finite size, subjective uncertainty remains when predicting transitions. 
It is precisely the subjective epistemic uncertainty which Bayesian modeling excels at, which helps mitigate overfitting.
Below, we describe how we use combinations of `probabilistic networks' to capture aleatoric uncertainty and `ensembles' to capture epistemic uncertainty. 
Each combination is summarized in \tab{table:modelclasses}.
		

\paragraph{Probabilistic neural networks (P)}
\label{sec:probabilistic}
We define a \textit{probabilistic} NN as a network whose output neurons simply parameterize a probability distribution function, capturing aleatoric uncertainty, and should not be confused with Bayesian inference. We use the negative log prediction probability as our loss function $\text{loss}_{\text{P}}(\parameters) = -\sum_{n=1}^{N} \log \dynamicsModel_{\parameters}(\MDPState_{n+\!1} | \MDPState_n,\MDPAction_n)$.
For example, we might define our predictive model to output a Gaussian distribution with diagonal covariances parameterized by $\parameters$ and conditioned on $\MDPState_n$ and $\MDPAction_n$, i.e.: \smash{$\dynamicsModel = \prob(\MDPState_{\MDPTime + 1} | \MDPState_{\MDPTime}, \MDPAction_{\MDPTime}) = \gauss{\vec{\mean}_{\theta}(\MDPState_{\MDPTime}, \MDPAction_{\MDPTime})}{\varMat_{\theta}(\MDPState_{\MDPTime}, \MDPAction_{\MDPTime})}$}.
Then the loss becomes 
%
\vspace{-5pt}
\begin{align}
\text{loss}_{\text{Gauss}}(\parameters) \!=\!\! \sum_{n=1}^{N} \left[\mean_{\parameters}(\MDPState_n,\!\MDPAction_n)\!-\! \MDPState_{n+\!1}\right]^\top\!\!\varMat^{-1}_{\parameters}(\MDPState_n,\!\MDPAction_n)\!\left[\mean_{\parameters}(\MDPState_n,\!\MDPAction_n)\!-\!\MDPState_{n+\!1}\right] \!+\! \log\det{\varMat_{\parameters}(\MDPState_n,\!\MDPAction_n)}.
% \vspace{-10pt}
\end{align}
%
Such network outputs, which in our particular case parameterizes a Gaussian distribution, models aleatoric uncertainty, otherwise known as heteroscedastic noise (meaning the output distribution is a function of the input). However, it does not model epistemic uncertainty, which cannot be captured with purely discriminative training.
Choosing a Gaussian distribution is a common choice for continuous-valued states, and reasonable if we assume that any stochasticity in the system is unimodal. 
However, in general, any tractable distribution class can be used. 
To provide for an expressive dynamics model, we can represent the parameters of this distribution (e.g., the mean and covariance of a Gaussian) as nonlinear, parametric functions of the current state and action, which can be arbitrarily complex but deterministic. 
This makes it feasible to incorporate NNs into a probabilistic dynamics model even for high-dimensional and continuous states and actions.
Finally, an under-appreciated detail of probabilistic networks is that their variance has \textit{arbitrary} values for out-of-distribution inputs, which can disrupt planning.
We discuss how to mitigate this issue in \app{sec:variance-bounding}.
%
    
\paragraph{Deterministic neural networks (D)} 
\label{sec:deterministic}
For comparison, we define a deterministic NN as a special-case probabilistic network that outputs delta distributions centered around point predictions denoted as \smash{$\dynamicsModel_{\parameters}(\MDPState_{\MDPTime}, \MDPAction_{\MDPTime})$}:
\smash{$\dynamicsModel_{\parameters}(\MDPState_{\MDPTime + 1} | \MDPState_{\MDPTime}, \MDPAction_{\MDPTime}) = \prob(\MDPState_{\MDPTime + 1} | \MDPState_{\MDPTime}, \MDPAction_{\MDPTime}) = \delta(\MDPState_{\MDPTime + 1} - \dynamicsModel_{\parameters}(\MDPState_{\MDPTime}, \MDPAction_{\MDPTime}))$}, trained using the MSE loss:
\smash{$ \text{loss}_{D}(\parameters) = \sum_{n=1}^{N}{\lVert \MDPState_{n+1} - \dynamicsModel_{\parameters}(\MDPState_n, \MDPAction_n) \rVert}$}.
%
Although MSE can be interpreted as $\text{loss}_{\text{P}}(\parameters)$ with a Gaussian model of fixed unit variance, in practice this variance cannot be used for uncertainty-aware propagation, since it does not correspond to any notion of uncertainty (e.g., a deterministic model with infinite data would be adding variance to particles for no good reason).

\paragraph{Ensembles (DE and PE)}
A principled means to capture epistemic uncertainty is with Bayesian inference. 
Whilst accurate Bayesian NN inference is possible with sufficient compute~\citep{neal1995bayesian}, approximate inference methods \citep{Blundell2015, gal2017concrete, hernandez2015probabilistic} have enjoyed recent popularity given their simpler implementation and faster training times.
Ensembles of bootstrapped models are even simpler still: given a base model, no additional (hyper-)parameters need be tuned, whilst still providing reasonable uncertainty estimates~\citep{efron1994introduction, Osband2016a,kurutach2018model}.
We consider ensembles of $\Bootstrap$-many bootstrap models, using \smash{$\parameters_\bootstrap$} to refer to the parameters of our \smash{$\bootstrap^{\text{th}}$} model \smash{$\dynamicsModel_{\parameters_\bootstrap}$}.
Ensembles can be composed of deterministic models (DE) or probabilistic models (PE) -- as done by \citet{Lakshminarayanan2017} -- both of which define predictive probability distributions: 
\smash{$\dynamicsModel_{\parameters} = \tfrac{1}{\Bootstrap}\sum_{\bootstrap=1}^\Bootstrap \dynamicsModel_{\parameters_\bootstrap}$}.
A visual example is provided in \app{app:pe-model-toy}.
Each of our bootstrap models have their unique dataset $\D_{\bootstrap}$, generated by sampling (with replacement) $N$ times the dynamics dataset recorded so far $\D$, where $N$ is the size of $\D$. We found $\Bootstrap=5$ sufficient for all our experiments.
To validate the number of layers and neurons of our models, we can visualize one-step predictions (e.g. \app{app:one-step-prediction}).
        
